\SourceLeaf{248}%
\Needspace{7\baselineskip}
\begin{center}
[\textsc{Gerling} \emph{an} \textsc{Gauss}.\quad \emph{Marburg, 26. Februar 1844.}]
\end{center}

\{\mbox{.\ .\ .\ .\ .\ .\ .\ .\ .} Seit ich 1813 auf die symmetrischen Dreiecke und die Fläche des sphärischen Dreiecks kam, hat mir immer eine scharfe Definition oder ein scharfes Kriterium der Symmetrie im Unterschied zur Kongruenz gefehlt. Seit letztem Winter bin ich nun darauf gekommen, die Sache auch im \textls{Vortrag} auf folgende Weise darzustellen. Für mich selbst war ich schon früher bei anderer Gelegenheit darauf gekommen (1834).

Jedes Gebilde in der Ebene beziehen wir auf \textls{zwei} Dimensionen (rechtwinklige Koordinaten), jedes räumliche auf \textls{drei}. Nehmen wir nun an, dass die Koordinaten ihr Vorzeichen wechseln können, so wird im Allgemeinen kongruent, was in \textls{zwei} Dimensionen das Vorzeichen wechselt. Symmetrisch wird dagegen, was in \textls{einer} oder in \textls{drei} Dimensionen das Vorzeichen wechselt. --- Hieraus folgt als Korollar,

1) dass bei Koordinaten von einem Punkt aus die Umkehr aller radii vectores (sit venia verbo) bei räumlichen Konstruktionen Symmetrie, bei ebenen aber Kongruenz ergeben muss,

\Needspace{7\baselineskip}
2) dass die Symmetrie in die Kongruenz übergeht, wenn sich das Gebilde durch eine Gerade oder eine Ebene in zwei symmetrische Hälften teilen lässt. ---

Sie würden mich sehr verbinden, wenn Sie mir gelegentlich Ihr Urteil über diese Darstellungsweise mitteilten.\}

\par\bigskip

\Needspace{7\baselineskip}
\begin{center}
\textsc{Gauss} an \textsc{Gerling}.\quad Göttingen, 8. April 1844.
\end{center}

\mbox{.\ .\ .\ .\ .\ .\ .\ .} Ihre Bemerkungen über Symmetrie und Kongruenz sind vollkommen zutreffend. Zu wünschen wäre noch der metaphysische Grund, \textls{warum} es so ist (was bei Ihnen nur als beobachtete Tatsache erscheint), und damit auch die Erweiterung auf eine Geometrie von mehr als 3 Dimensionen, für die wir Menschen keine Anschauung haben, die aber\space%
